\section{Correspondência entre os três espaços}
A arquitetura distingue a variedade de ondas $\mathcal W$, o espaço de estados
$\mathcal U\simeq\mathbb R^2_{(u,v)}$ e o espaço físico
$\mathcal X\simeq\mathbb R^2_{(x,t)}$.  Para uma órbita de rarefação
$\gamma\subset\mathcal C$, a projeção em estados é
\[
  \mathcal R=\bar U\circ\gamma,
\]
pois $Y=0$ implica $U_-=U_+=\bar U$.  Quando a velocidade característica
$s$ é monótona ao longo do trecho considerado, a realização auto-similar usa
\[
  \xi=\frac{x}{t},\qquad \xi=s(U),\qquad t>0.
\]
Assim, o mesmo objeto é representado sucessivamente em
\[
  \mathcal W\longrightarrow\mathcal U\longrightarrow\mathcal X.
\]
A implementação desta correspondência não constitui, por si só, um solver do
problema de Riemann geral. Órbitas com reversão de $s$ são rejeitadas como
leques físicos regulares até que uma construção matemática apropriada seja
especificada.
